% Part: computability
% Chapter: computability-theory
% Section: k-1

\documentclass[../../../include/open-logic-section]{subfiles}

\begin{document}

\olfileid[pt]{cmp}{thy}{k1}
\olsection{Um exemplo de redutibilidade}

Consideremos uma aplicação da \olref[ppr]{prop:reduce}.

\begin{prop}
\ollabel{prop:k1}
Seja
\[
K_1 = \Setabs{e}{\cfind{e}(0) \fdefined}.
\]
Então $K_1$ é computavelmente enumerável, mas não é computável.
\end{prop}

\begin{proof}
Como $K_1 = \Setabs{e}{\lexists[s][T(e,0,s)]}$, $K_1$ é computavelmente
enumerável por \olref[eqc]{thm:exists-char}.

Para mostrar que $K_1$ não é computável, mostraremos que $K_0$ é
redutível a ele.

\begin{explain}
Isso é um pouco complicado, pois, usando $K_1$, só podemos fazer
perguntas sobre computações que começam com uma entrada específica,
$0$. Suponha que você tenha um amigo inteligente capaz de responder
a perguntas desse tipo (amigos assim são conhecidos como
``oráculos''). Suponha, então, que alguém se aproxime e lhe pergunte
se $\tuple{e, x}$ pertence ou não a $K_0$, isto é, se a máquina $e$
para ou não na entrada $x$. Uma coisa que você pode fazer é construir
outra máquina, $e_x$, que, para \emph{qualquer} entrada, ignore essa
entrada e, em vez disso, execute~$e$ na entrada $x$. Então, é claro,
a pergunta sobre se a máquina $e$ para na entrada $x$ equivale à
pergunta sobre se a máquina $e_x$ para na entrada $0$ (ou em qualquer
outra entrada). Você pergunta então a seu amigo se essa nova máquina,
$e_x$, para na entrada $0$; a resposta de seu amigo à pergunta
modificada fornece a resposta à pergunta original. Isso fornece a
redução desejada de $K_0$ a~$K_1$.
\end{explain}

Usando a função parcialmente computável universal, seja $f$
a função ternária definida por
\[
f(x,y,z) \simeq \cfind{x}(y).
\]
Observe que $f$ ignora completamente sua terceira entrada. Escolha um
índice $e$ tal que $f = \cfind{e}[3]$; assim, temos
\[
\cfind{e}[3](x,y,z) \simeq \cfind{x}(y).
\]
Pelo teorema $s$-$m$-$n$, existe uma função $s(e,x,y)$ tal que, para
todo $z$,
\begin{align*}
\cfind{s(e,x,y)}(z) & \simeq \cfind{e}[3](x,y,z) \\
& \simeq \cfind{x}(y).
\end{align*}

\begin{explain}
Em termos do argumento informal acima, $s(e,x,y)$ é um índice da
máquina que, para qualquer entrada $z$, ignora essa entrada e calcula
$\cfind{x}(y)$.
\end{explain}

Em particular, temos
\[
\cfind{s(e,x,y)}(0) \fdefined \quad \text{se, e somente se,} \quad
\cfind{x}(y) \fdefined.
\]
Em outras palavras, $\tuple{x, y} \in K_0$ se, e somente se,
$s(e,x,y) \in K_1$. Logo, a função $g$ definida por
\[
g(w) = s(e,(w)_0,(w)_1)
\]
é uma redução de $K_0$ a~$K_1$.
\end{proof}

\end{document}

